%! The Complete Solution to Ax = b — Unit 3, Lesson 3.3 — Homework
\documentclass[10pt]{article}
\usepackage{linalg-article}
\usepackage{linalg-boxes}
\newcommand{\TallMath}[1]{$\displaystyle #1\rule[-1.4em]{0pt}{3.2em}$}
\newcommand{\xx}{\mathbf{x}}
\newcommand{\yy}{\mathbf{y}}
\newcommand{\bb}{\mathbf{b}}
\newcommand{\svec}{\mathbf{s}}
\newcommand{\zero}{\mathbf{0}}

\begin{document}

\pageheader{Unit 3, Lesson 3.3}{Homework}
\namedateperiod

\vspace{0.12in}

\begin{notesbox}{Practice}
Reduce $[A\mid\bb]$, find $\xx_p$ (free variables $=0$), attach the nullspace, and \emph{justify}. Always
say which shape (point, line, or plane) the solution set is, and whether it passes through the origin.

\begin{enumerate}[leftmargin=1.9em, itemsep=11pt]
  \item \textbf{Complete solutions.} For each system, give $\xx_p$, the nullspace part, the complete
        solution, and its shape.
    \begin{enumerate}[label=(\alph*), leftmargin=1.8em, itemsep=6pt]
      \item $\begin{bmatrix} 1 & 2 \\ 1 & 3 \end{bmatrix}\xx=\begin{bmatrix} 3 \\ 4 \end{bmatrix}$ \vspace{0.9cm}
      \item $\begin{bmatrix} 1 & 2 \\ 2 & 4 \end{bmatrix}\xx=\begin{bmatrix} 2 \\ 4 \end{bmatrix}$ \vspace{0.9cm}
      \item $\begin{bmatrix} 1 & 2 \\ 2 & 4 \end{bmatrix}\xx=\begin{bmatrix} 2 \\ 5 \end{bmatrix}$ \vspace{0.9cm}
      \item $\begin{bmatrix} 1 & 1 & 2 \\ 0 & 1 & 1 \end{bmatrix}\xx=\begin{bmatrix} 4 \\ 1 \end{bmatrix}$ \vspace{1.1cm}
    \end{enumerate}
  \item \textbf{Count without solving.} The system $A\xx=\bb$ is solvable. State how many solutions.
    \begin{enumerate}[label=(\alph*), leftmargin=1.8em, itemsep=4pt]
      \item $A$ has $3$ columns and rank $r=3$. \hfill
      \item $A$ has $5$ columns and rank $r=3$. \hfill
    \end{enumerate}
    \vspace{0.7cm}
  \item \textbf{Solvable or not?} For $A=\begin{bmatrix} 1 & 3 \\ 2 & 6 \end{bmatrix}$, decide which of
        $\bb=\begin{bsmallmatrix} 2\\4 \end{bsmallmatrix}$ and $\bb=\begin{bsmallmatrix} 2\\5 \end{bsmallmatrix}$
        gives a solvable system. Reduce $[A\mid\bb]$ in each case to justify. \writelines{2}
  \item \textbf{Explain.} When $\bb\ne\zero$, why is the solution set of $A\xx=\bb$ \emph{never} a
        subspace, even though the nullspace always is? What is the one thing it is missing? \writelines{2}
\end{enumerate}
\end{notesbox}

\begin{extensionbox}
\textbf{Rank and the number of solutions.}
\begin{enumerate}[label=(\alph*), leftmargin=1.8em, itemsep=5pt]
  \item \textbf{Four cases.} For a solvable $A\xx=\bb$, say how many solutions each rank situation forces
        (one, or infinitely many): (i) full column rank $r=n$; (ii) $r<n$. Then say which case
        \emph{always} has a solution for every $\bb$: full row rank $r=m$.
  \item \textbf{Two means infinitely many.} Show that if $A\xx=\bb$ has \emph{two different} solutions
        $\xx_1\ne\xx_2$, then it has infinitely many. (Hint: what is $\xx_1-\xx_2$?)
\end{enumerate}
\writelines{3}
\end{extensionbox}

\begin{spiralbox}
\textbf{Coming next --- Lesson 3.4: Independence, Basis, and Dimension.} Today every solution was
$\xx_p$ plus a combination of the special solutions. Next: which of those special solutions are
\emph{really} independent (none a combination of the others), and the smallest set that still spans a
space --- a \emph{basis} --- with the count called the \emph{dimension}.
\end{spiralbox}

\end{document}
